%% notes-tex/027-metabolic-flux-module/sections/4-arms.tex
%% [[experiments.027-betaxanthin-metabolic]]
%% https://github.com/Mjvolk3/torchcell/tree/main/notes-tex/027-metabolic-flux-module/sections/4-arms.tex

\section{The ablation registry\secstatus{todo}}
\label{sec:arms}

\subsection{What each arm turns on\secstatus{todo}}

The registry is one dictionary, \texttt{ARMS} in
\texttt{experiments/026-metabolism-flux/scripts/train\_flux.py}, and it is the single
source of truth for what the layer does in any run. An experiment names an arm; it never
redefines one. That is what keeps a difference between two arms attributable to the module
rather than to a hyperparameter that moved alongside it.

\begin{table}[H]
  \centering
  \footnotesize
  \caption[The arms]{The registry as implemented. Each row is a strict addition to the one
  above it except \texttt{nullspace}, which changes the parameterization rather than the
  physics.}
  \label{tab:arms}
  %% SOURCE: experiments/026-metabolism-flux/scripts/train_flux.py, ARMS.
  \begin{tabular}{@{}lllll@{}}
    \hline
    Arm & Flux & Thermo & Capacity, budget, dissipation & Isolates \\
    \hline
    \texttt{pooled} & off & -- & -- & the no-module control \\
    \texttt{flux\_off} & on & \texttt{OFF} & off & routing through $S$ alone \\
    \texttt{flux\_free} & on & \texttt{FREE} & off & structural loop-freedom \\
    \texttt{flux\_anchored} & on & \texttt{ANCHORED} & on & \textbf{the treatment} \\
    \texttt{flux\_nullspace} & on & \texttt{ANCHORED} & on & the exactness budget \\
    \hline
  \end{tabular}
\end{table}

The ladder is deliberate. \texttt{flux\_off} routes the representation through the
stoichiometric network with nothing but mass balance and the GPR, so it asks whether the
network's topology alone carries signal. \texttt{flux\_free} adds a learned potential,
which buys loop-freedom and no chemistry. \texttt{flux\_anchored} replaces the learned
potential with tabulated formation energies and switches on the enzyme capacity bound, the
protein budget and the dissipation limit. \texttt{flux\_nullspace} keeps the anchored
physics but spends the exactness budget on $S\nu = 0$ instead of on the box, which is the
one arm that tests the parameterization choice rather than the model.

\subsection{The control that is missing, and why it matters\secstatus{todo}}

Every arm above changes the module's \emph{physics} together with its \emph{parameter
count}. So a gap between \texttt{flux\_anchored} and \texttt{pooled} is consistent with two
different explanations: the tabulated chemistry is informative, or the extra structure acts
as a capacity constraint that would help with any numbers in it.

The arm that separates those is \texttt{flux\_shuffled}: the full anchored module with
$k_{cat}$ \textbf{permuted across reactions}, holding architecture, parameter count and
optimization fixed and destroying only the correspondence between a reaction and its
turnover number. If \texttt{flux\_anchored} and \texttt{flux\_shuffled} land together, the
kinetic tables are decoration.

It is \textbf{not in the registry and not in the current run.} Adding it is a code change
that was not worth making untested, and it is the first item of the next round.

\subsection{The permutation nulls\secstatus{todo}}

Two arms exist only to be compared against: \texttt{null\_pooled} and
\texttt{null\_anchored} train on \textbf{permuted labels} while validating on real ones.
They matter because of what 026 measured. Against zero, its flux-versus-pooled gap of
$+0.0735$ is 7.7 standard errors. Against its own permuted-label gap of $+0.0302$, the
excess is $+0.0433$ at 1.9 standard errors. More than 40 percent of ``the flux module
helps'' survived destroying the labels.

So the decision rule is a difference of differences,
\begin{equation}
  \big(\text{flux} - \text{pooled}\big) - \big(\text{null flux} - \text{null pooled}\big),
\end{equation}
and a gap measured against zero is not evidence.
